\section{Discussion}\label{sec:discussion}

\subsection{What the Study Establishes}
\paragraph*{Empty difference sets carry substantial false-assurance risk} On the development corpus, about one in six plain no-difference outcomes concealed a labelled behavioural change, and about one in five of the V7-B gate's reassuring verdicts did. These rates depend on construction labels and on a corpus dominated by small generated programs, so their magnitude does not transfer. On these data the rate was neither negligible nor demonstrably reduced by stricter abstention; whether this holds elsewhere is untested. A pipeline that treats an empty difference set as a pass is accepting a risk that it does not measure.

\paragraph*{Abstention is not reliability} The V7-B gate issued a third as many \NDO{} verdicts as plain differential testing without lowering the risk among those it issued, and full \cec{} at the same coverage did not beat random abstention (observed risk 4.4 points higher, not significant). Had we compared these procedures at their own operating points, the gate would have appeared to reduce the absolute number of false assurances from 31 to 13, which is true and almost entirely explained by issuing fewer verdicts. Matched-coverage comparison removed that illusion. We think this methodological point generalises beyond \sbg{}: any proposal that withholds verdicts---through confidence thresholds, sufficiency gates or LLM-based filters---should report risk at matched coverage against random abstention.

\paragraph*{Evidence-response properties are useful as conformance tests} The \ers{} matrix found a real defect in the evidence profile that no aggregate metric revealed, distinguished a stored verdict from a recomputed one and both from the \sbg{} trust gate, and made the dependence of the fitted model on a single feature visible. These are label-free checks that any estimator of verdict reliability can be subjected to. In our view this is the most transferable contribution of the paper: a short list of properties that a no-divergence risk estimator should satisfy, a concrete perturbation protocol that tests them, and evidence that the protocol finds defects.

\subsection{What the Study Does Not Establish}
The study does not show that change-conditioned evidence calibration predicts false assurance better than coverage alone, and it does not show that \ers{} features add predictive value. Both effects were estimated with intervals that include zero and with seed-to-seed sign changes. With 61 issued verdicts the study could only have detected differences of roughly 16 points, so it is uninformative about effects of the 5-point size that would matter in practice. The placebo comparison hints that the \ers{} features are not pure noise, but the evidence is weak and unadjusted. The appropriate reading is that the predictive hypotheses remain open.

Several design facts also limit the predictive comparison. The fitted model retains essentially one feature, multiplicity; the provenance, integrity and environment features were constant in the training data and cannot have contributed; and the observable-influence feature is a reachability proxy rather than a dependence analysis. A richer, more diverse corpus is a precondition for testing whether the full profile matters.

\subsection{Implications}
\paragraph*{For practitioners} A no-difference result should be reported together with the evidence that supports it: at minimum, whether the changed code was executed, by how many observably distinct inputs, and on which channels. Multiplicity---the number of distinct observations reaching each changed unit---was the one feature that fitted models consistently retained, and it is cheap to compute from line-level tracing, although our data do not show that it lowers matched-coverage risk. Integrity failures should block reassuring verdicts outright rather than being folded into a score.

\paragraph*{For researchers} Tools that generate regression tests, including recent LLM-based tools~\cite{pradel2026testora,hu2026difftestgen,ayub2026semadiff}, are evaluated almost exclusively on the divergences they find. Their behaviour on changes where they find nothing deserves equal attention, and the estimand of Section~\ref{sec:problem} offers a way to measure it. Their generated tests could also serve as evidence for \sbg{}, which would give the profile more diverse inputs than the declared tests used here.

\subsection{Toward a Confirmatory Study}
The power analysis sizes the confirmatory study at several hundred to over a thousand issued \NDO{} verdicts, adjudicated independently. The infrastructure for it exists---the blinded two-adjudicator protocol, the prediction lock and the evaluation code---but four prerequisites must be met first: (i)~an external corpus of real benchmark checkouts executed by the collector, rather than extracted units; (ii)~an adjudication package from which labels cannot be inferred; (iii)~a prediction lock with an external timestamp, created before any adjudication information exists; and (iv)~a profile in which provenance and environment are populated from the runs, and an observable-influence feature based on dynamic dependence.
